% ============================================================
% arXiv & Daily Papers daily note template
% ============================================================

\section{arXiv \& Daily Papers 日报：YYYY年MM月DD日}

\begin{conceptbox}
\textbf{方向}：后训练 · 视频生成/编辑 · 图像生成/编辑 · 统一理解与生成 · 世界模型 · 具身智能 · Style\\
\textbf{数据来源}：\href{https://arxiv.org/list/cs.CV/new}{arXiv cs.CV} · \href{https://arxiv.org/list/cs.CL/new}{cs.CL} · \href{https://arxiv.org/list/cs.AI/new}{cs.AI} · \href{https://arxiv.org/list/cs.RO/new}{cs.RO} · \href{https://arxiv.org/list/cs.LG/new}{cs.LG} · \href{https://huggingface.co/papers}{Hugging Face Daily Papers}\\
\textbf{整理日期}：YYYY-MM-DD
\end{conceptbox}

\subsection{今日重点}

\begin{enumerate}[leftmargin=1.5em]
  \item \textbf{论文名}：一句话说明为什么重要。\\
  \textbf{arXiv}：\href{https://arxiv.org/abs/xxxx.xxxxx}{xxxx.xxxxx}
\end{enumerate}

\subsection{按方向归类}

\subsubsection{后训练}

\begin{itemize}[leftmargin=1.5em]
  \item \textbf{论文名}：核心贡献。
\end{itemize}

\subsubsection{视频生成/编辑}

\begin{itemize}[leftmargin=1.5em]
  \item \textbf{论文名}：核心贡献。
\end{itemize}

\subsubsection{图像生成/编辑与统一理解生成}

\begin{itemize}[leftmargin=1.5em]
  \item \textbf{论文名}：核心贡献。
\end{itemize}

\subsubsection{世界模型与具身智能}

\begin{itemize}[leftmargin=1.5em]
  \item \textbf{论文名}：核心贡献。
\end{itemize}

\subsubsection{Style / Voice / Identity}

\begin{itemize}[leftmargin=1.5em]
  \item \textbf{论文名}：核心贡献。
\end{itemize}

\subsection{值得深读}

\begin{itemize}[leftmargin=1.5em]
  \item \textbf{优先级 A}：
  \item \textbf{优先级 B}：
\end{itemize}

\subsection{暂时略过}

\begin{itemize}[leftmargin=1.5em]
  \item 与当前主线弱相关，先记录标题即可。
\end{itemize}

\subsection{后续动作}

\begin{itemize}[leftmargin=1.5em]
  \item 是否需要单篇深读。
  \item 是否需要补到某个专题学习 session。
\end{itemize}
